\section{Exercices d'application}

\rubrique{Vocabulaire et représentations graphiques}

\exo{}\diff{1}
On étudie le nombre d'enfants par famille dans un quartier de Yaoundé. L'enquête porte sur 30 familles. Les résultats sont les suivants :
\[
2; 1; 3; 0; 2; 4; 1; 2; 3; 1; 2; 0; 2; 1; 3; 2; 1; 2; 0; 3; 2; 1; 4; 2; 1; 2; 3; 2; 1; 2
\]
\begin{enumerate}[label=\alph*)]
\item Déterminer la population étudiée, le caractère et sa nature.
\item Dresser le tableau des effectifs et des fréquences (en pourcentage).
\item Représenter la série par un diagramme en bâtons.
\end{enumerate}

\exo{}\diff{1}
Dans une classe de Première C du lycée de Bafoussam, on a relevé les notes obtenues par les élèves à un devoir de mathématiques :
\[
8; 12; 15; 7; 10; 14; 9; 11; 13; 6; 16; 10; 8; 12; 14; 9; 11; 13; 10; 7; 15; 8; 12; 10; 11
\]
\begin{enumerate}[label=\alph*)]
\item Regrouper ces données en classes d'amplitude 2 (de 6 à 8, 8 à 10, etc.).
\item Construire l'histogramme et le polygone des effectifs.
\end{enumerate}

\rubrique{Paramètres de position}

\exo{}\diff{1}
On considère la série statistique suivante :
\[
x_i : 3; 5; 7; 9; 11 \quad \text{avec effectifs } n_i : 2; 5; 8; 4; 1
\]
Calculer la moyenne arithmétique de cette série.

\exo{}\diff{2}
Les salaires mensuels (en milliers de francs CFA) des employés d'une entreprise de Douala sont donnés par le tableau suivant :
\begin{center}
\begin{tabular}{|c|c|c|c|c|c|}
\hline
Salaire & [50; 100[ & [100; 150[ & [150; 200[ & [200; 250[ & [250; 300[ \\
\hline
Effectif & 12 & 25 & 30 & 18 & 5 \\
\hline
\end{tabular}
\end{center}
\begin{enumerate}[label=\alph*)]
\item Calculer le salaire moyen.
\item Déterminer la classe modale.
\end{enumerate}

\exo{}\diff{2}
On donne la série ordonnée des notes de 15 élèves :
\[
5; 6; 8; 9; 10; 10; 11; 12; 13; 14; 14; 15; 16; 17; 18
\]
\begin{enumerate}[label=\alph*)]
\item Déterminer la médiane de cette série.
\item Calculer le premier quartile \(Q_1\) et le troisième quartile \(Q_3\).
\end{enumerate}

\rubrique{Paramètres de dispersion}

\exo{}\diff{2}
Pour la série de l'exercice 3, calculer :
\begin{enumerate}[label=\alph*)]
\item L'étendue.
\item La variance et l'écart-type.
\end{enumerate}

\exo{}\diff{3}
Deux groupes d'élèves ont passé un test de mathématiques. Les résultats sont :
\begin{itemize}
\item Groupe A : moyenne = 12, écart-type = 2,5
\item Groupe B : moyenne = 12, écart-type = 1,8
\end{itemize}
\begin{enumerate}[label=\alph*)]
\item Quel groupe est le plus homogène ? Justifier.
\item Si un élève du groupe A a 14 et un élève du groupe B a 13, lequel est le mieux classé par rapport à son groupe ? (On pourra calculer les notes centrées réduites.)
\end{enumerate}

\exo{}\diff{3}
On considère la série statistique suivante :
\[
x_i : 2; 5; 8; 11; 14 \quad \text{avec effectifs } n_i : 3; 7; 10; 6; 4
\]
\begin{enumerate}[label=\alph*)]
\item Calculer la moyenne \(\bar{x}\).
\item Calculer la variance et l'écart-type.
\item Déterminer l'intervalle \([\bar{x} - \sigma; \bar{x} + \sigma]\) et le pourcentage de valeurs dans cet intervalle.
\end{enumerate}

\section{Exercices d'entraînement}

\rubrique{Vocabulaire et représentations}

\exo{}\diff{1}
Une enquête sur le nombre de livres lus par mois par 40 élèves d'un lycée de Garoua donne les résultats suivants :
\[
0; 1; 2; 1; 3; 0; 2; 1; 1; 2; 3; 1; 0; 2; 1; 2; 1; 0; 3; 1; 2; 1; 1; 2; 0; 1; 2; 3; 1; 2; 1; 0; 2; 1; 1; 2; 3; 1; 2; 1
\]
\begin{enumerate}[label=\alph*)]
\item Identifier la population, le caractère et sa nature.
\item Construire le tableau des effectifs et des fréquences.
\item Représenter par un diagramme circulaire.
\end{enumerate}

\exo{}\diff{1}
Les tailles (en cm) de 50 élèves d'une classe de Première C sont :
\[
152; 158; 160; 155; 162; 157; 153; 159; 161; 156; 154; 160; 158; 155; 163; 157; 152; 159; 161; 156; 155; 158; 160; 154; 162; 157; 153; 159; 161; 155; 158; 156; 160; 154; 162; 157; 153; 159; 161; 155; 158; 156; 160; 154; 162; 157; 153; 159; 161; 155
\]
\begin{enumerate}[label=\alph*)]
\item Regrouper en classes d'amplitude 3 (de 150 à 153, 153 à 156, etc.).
\item Tracer l'histogramme et le polygone des effectifs.
\end{enumerate}

\exo{}\diff{2}
Dans une entreprise de Yaoundé, on a relevé les âges des employés :
\begin{center}
\begin{tabular}{|c|c|c|c|c|c|}
\hline
Âge & [20; 25[ & [25; 30[ & [30; 35[ & [35; 40[ & [40; 45[ \\
\hline
Effectif & 8 & 15 & 20 & 12 & 5 \\
\hline
\end{tabular}
\end{center}
Représenter cette série par un histogramme et un polygone des effectifs cumulés croissants.

\rubrique{Paramètres de position}

\exo{}\diff{1}
Calculer la moyenne de la série : \(x_i : 4; 7; 10; 13\) avec effectifs \(n_i : 3; 6; 9; 2\).

\exo{}\diff{2}
Les notes d'un devoir sont :
\[
6; 8; 9; 10; 12; 13; 14; 15; 16; 17; 18; 19; 20
\]
Déterminer la médiane, \(Q_1\) et \(Q_3\).

\exo{}\diff{2}
Une série statistique a pour moyenne 15. Si on ajoute 3 à chaque valeur, quelle est la nouvelle moyenne ? Si on multiplie chaque valeur par 2, quelle est la nouvelle moyenne ?

\exo{}\diff{2}
On donne la série :
\[
x_i : 1; 2; 3; 4; 5; 6 \quad \text{avec effectifs } n_i : 2; 5; 8; 6; 3; 1
\]
Déterminer le mode et la médiane.

\exo{}\diff{3}
Dans une usine de Limbé, les salaires (en milliers de francs CFA) sont :
\begin{center}
\begin{tabular}{|c|c|c|c|c|}
\hline
Salaire & [80; 120[ & [120; 160[ & [160; 200[ & [200; 240[ \\
\hline
Effectif & 10 & 25 & 15 & 5 \\
\hline
\end{tabular}
\end{center}
Calculer le salaire moyen et la classe modale.

\rubrique{Paramètres de dispersion}

\exo{}\diff{2}
Pour la série : \(x_i : 2; 4; 6; 8; 10\) avec effectifs \(n_i : 1; 3; 5; 3; 1\), calculer l'étendue, la variance et l'écart-type.

\exo{}\diff{2}
Deux séries A et B ont pour caractéristiques :
\begin{itemize}
\item Série A : moyenne = 10, variance = 4
\item Série B : moyenne = 10, variance = 9
\end{itemize}
Laquelle est la plus dispersée ? Justifier.

\exo{}\diff{3}
On considère la série :
\[
x_i : 0; 3; 6; 9; 12 \quad \text{avec effectifs } n_i : 4; 8; 12; 6; 2
\]
\begin{enumerate}[label=\alph*)]
\item Calculer la moyenne et l'écart-type.
\item Déterminer le pourcentage de valeurs dans l'intervalle \([\bar{x} - \sigma; \bar{x} + \sigma]\).
\end{enumerate}

\exo{}\diff{3}
Les notes de deux classes sont :
\begin{itemize}
\item Classe A : moyenne = 11, écart-type = 2
\item Classe B : moyenne = 11, écart-type = 3
\end{itemize}
Un élève de A a 13 et un élève de B a 14. Lequel est le meilleur relativement à sa classe ?

\section{Problèmes type Probatoire/Bac}

\sujetbac[Statistiques et interprétation]
Dans une entreprise de transport à Douala, on a relevé les distances parcourues (en km) par 50 camions en une journée :
\begin{center}
\begin{tabular}{|c|c|c|c|c|c|}
\hline
Distance & [50; 100[ & [100; 150[ & [150; 200[ & [200; 250[ & [250; 300[ \\
\hline
Effectif & 5 & 12 & 18 & 10 & 5 \\
\hline
\end{tabular}
\end{center}
\begin{enumerate}[label=\arabic*)]
\item Représenter cette série par un histogramme.
\item Calculer la distance moyenne parcourue.
\item Déterminer la classe modale.
\item Calculer la variance et l'écart-type.
\item Quel est le pourcentage de camions parcourant une distance comprise entre \(\bar{x} - \sigma\) et \(\bar{x} + \sigma\) ?
\end{enumerate}

\sujetbac[Comparaison de séries]
Dans deux lycées de Yaoundé, on a relevé les notes d'un test de mathématiques :
\begin{itemize}
\item Lycée A : moyenne = 12, écart-type = 2,5
\item Lycée B : moyenne = 11, écart-type = 1,8
\end{itemize}
\begin{enumerate}[label=\arabic*)]
\item Quel lycée a la meilleure moyenne ?
\item Quel lycée est le plus homogène ? Justifier.
\item Un élève du lycée A a 15 et un élève du lycée B a 14. Calculer leur note centrée réduite. Lequel est le mieux classé par rapport à son lycée ?
\item Si on ajoute 2 points à chaque note du lycée A, que deviennent la moyenne et l'écart-type ?
\end{enumerate}

\sujetbac[Statistiques et décision]
Un gérant de supermarché à Bafoussam étudie le montant des achats (en milliers de francs CFA) de 100 clients :
\begin{center}
\begin{tabular}{|c|c|c|c|c|c|}
\hline
Montant & [0; 10[ & [10; 20[ & [20; 30[ & [30; 40[ & [40; 50[ \\
\hline
Effectif & 15 & 25 & 35 & 20 & 5 \\
\hline
\end{tabular}
\end{center}
\begin{enumerate}[label=\arabic*)]
\item Calculer le montant moyen des achats.
\item Déterminer la médiane et les quartiles.
\item Le gérant souhaite offrir une réduction aux 25\% des clients qui dépensent le plus. À partir de quel montant un client bénéficie-t-il de cette réduction ?
\item Calculer l'écart-type et interpréter le résultat.
\end{enumerate}

\section{Corrigés}

\corrige{de l'exercice 1}
\begin{enumerate}[label=\alph*)]
\item Population : les 30 familles du quartier. Caractère : nombre d'enfants par famille. Nature : quantitatif discret.
\item Tableau :
\begin{center}
\begin{tabular}{|c|c|c|c|c|c|}
\hline
Nombre d'enfants & 0 & 1 & 2 & 3 & 4 \\
\hline
Effectif & 3 & 9 & 12 & 4 & 2 \\
\hline
Fréquence (\%) & 10 & 30 & 40 & 13,33 & 6,67 \\
\hline
\end{tabular}
\end{center}
\item Diagramme en bâtons : axe horizontal pour les valeurs 0 à 4, axe vertical pour les effectifs, avec des bâtons de hauteurs respectives 3, 9, 12, 4, 2.
\end{enumerate}

\corrige{de l'exercice 2}
\begin{enumerate}[label=\alph*)]
\item Classes : [6; 8[, [8; 10[, [10; 12[, [12; 14[, [14; 16[. Effectifs : 2; 5; 8; 6; 4.
\item Histogramme : rectangles de largeur 2 et hauteurs 2, 5, 8, 6, 4. Polygone : points aux centres des classes (7; 2), (9; 5), (11; 8), (13; 6), (15; 4) reliés.
\end{enumerate}

\corrige{de l'exercice 3}
Moyenne : \(\bar{x} = \frac{3\times 2 + 5\times 5 + 7\times 8 + 9\times 4 + 11\times 1}{2+5+8+4+1} = \frac{6+25+56+36+11}{20} = \frac{134}{20} = 6,7\).

\corrige{de l'exercice 4}
\begin{enumerate}[label=\alph*)]
\item Centres des classes : 75; 125; 175; 225; 275. Moyenne : \(\bar{x} = \frac{75\times 12 + 125\times 25 + 175\times 30 + 225\times 18 + 275\times 5}{12+25+30+18+5} = \frac{900+3125+5250+4050+1375}{90} = \frac{14700}{90} \approx 163,33\) milliers de francs CFA.
\item Classe modale : [150; 200[ (effectif le plus élevé : 30).
\end{enumerate}

\corrige{de l'exercice 5}
\begin{enumerate}[label=\alph*)]
\item Effectif total \(N = 15\). Médiane : \( \frac{15+1}{2} = 8\)e valeur, soit 12.
\item \(Q_1\) : \(\frac{15}{4} = 3,75\) → 4e valeur : 9. \(Q_3\) : \(\frac{3\times 15}{4} = 11,25\) → 12e valeur : 15.
\end{enumerate}

\corrige{de l'exercice 6}
\begin{enumerate}[label=\alph*)]
\item Étendue : \(11 - 3 = 8\).
\item Variance : \(V = \frac{3^2\times 2 + 5^2\times 5 + 7^2\times 8 + 9^2\times 4 + 11^2\times 1}{20} - (6,7)^2 = \frac{18+125+392+324+121}{20} - 44,89 = \frac{980}{20} - 44,89 = 49 - 44,89 = 4,11\). Écart-type : \(\sigma = \sqrt{4,11} \approx 2,03\).
\end{enumerate}

\corrige{de l'exercice 7}
\begin{enumerate}[label=\alph*)]
\item Groupe B est plus homogène car son écart-type (1,8) est plus petit que celui de A (2,5).
\item Note centrée réduite pour A : \(z_A = \frac{14-12}{2,5} = 0,8\). Pour B : \(z_B = \frac{13-12}{1,8} \approx 0,56\). L'élève de A est mieux classé car son \(z\) est plus grand.
\end{enumerate}

\corrige{de l'exercice 8}
\begin{enumerate}[label=\alph*)]
\item Moyenne : \(\bar{x} = \frac{2\times 3 + 5\times 7 + 8\times 10 + 11\times 6 + 14\times 4}{3+7+10+6+4} = \frac{6+35+80+66+56}{30} = \frac{243}{30} = 8,1\).
\item Variance : \(V = \frac{2^2\times 3 + 5^2\times 7 + 8^2\times 10 + 11^2\times 6 + 14^2\times 4}{30} - (8,1)^2 = \frac{12+175+640+726+784}{30} - 65,61 = \frac{2337}{30} - 65,61 = 77,9 - 65,61 = 12,29\). Écart-type : \(\sigma = \sqrt{12,29} \approx 3,51\).
\item Intervalle : \([8,1 - 3,51; 8,1 + 3,51] = [4,59; 11,61]\). Valeurs dans cet intervalle : 5, 8, 11 (effectifs 7+10+6 = 23). Pourcentage : \(\frac{23}{30} \times 100 \approx 76,67\%\).
\end{enumerate}

\corrige{du problème 1}
\begin{enumerate}[label=\arabic*)]
\item Histogramme : rectangles de largeur 50 et hauteurs 5, 12, 18, 10, 5.
\item Centres : 75; 125; 175; 225; 275. Moyenne : \(\bar{x} = \frac{75\times 5 + 125\times 12 + 175\times 18 + 225\times 10 + 275\times 5}{50} = \frac{375+1500+3150+2250+1375}{50} = \frac{8650}{50} = 173\) km.
\item Classe modale : [150; 200[ (effectif 18).
\item Variance : \(V = \frac{75^2\times 5 + 125^2\times 12 + 175^2\times 18 + 225^2\times 10 + 275^2\times 5}{50} - 173^2 = \frac{28125+187500+551250+506250+378125}{50} - 29929 = \frac{1651250}{50} - 29929 = 33025 - 29929 = 3096\). Écart-type : \(\sigma = \sqrt{3096} \approx 55,64\) km.
\item Intervalle : \([173 - 55,64; 173 + 55,64] = [117,36; 228,64]\). Classes concernées : [100; 150[, [150; 200[, [200; 250[ (effectifs 12+18+10 = 40). Pourcentage : \(\frac{40}{50} \times 100 = 80\%\).
\end{enumerate}

\corrige{du problème 2}
\begin{enumerate}[label=\arabic*)]
\item Lycée A a la meilleure moyenne (12 > 11).
\item Lycée B est plus homogène (écart-type 1,8 < 2,5).
\item Pour A : \(z_A = \frac{15-12}{2,5} = 1,2\). Pour B : \(z_B = \frac{14-11}{1,8} \approx 1,67\). L'élève de B est mieux classé.
\item Nouvelle moyenne : \(12 + 2 = 14\). Nouvel écart-type : inchangé (2,5) car ajout d'une constante ne modifie pas la dispersion.
\end{enumerate}

\corrige{du problème 3}
\begin{enumerate}[label=\arabic*)]
\item Centres : 5; 15; 25; 35; 45. Moyenne : \(\bar{x} = \frac{5\times 15 + 15\times 25 + 25\times 35 + 35\times 20 + 45\times 5}{100} = \frac{75+375+875+700+225}{100} = \frac{2250}{100} = 22,5\) milliers de francs CFA.
\item Effectif total 100. Médiane : 50e et 51e valeurs dans [20; 30[. \(M_e = 20 + \frac{50-40}{35} \times 10 \approx 22,86\). \(Q_1\) : 25e valeur dans [10; 20[, \(Q_1 = 10 + \frac{25-15}{25} \times 10 = 14\). \(Q_3\) : 75e valeur dans [30; 40[, \(Q_3 = 30 + \frac{75-75}{20} \times 10 = 30\).
\item Les 25\% qui dépensent le plus sont au-dessus de \(Q_3\), soit à partir de 30 milliers de francs CFA.
\item Variance : \(V = \frac{5^2\times 15 + 15^2\times 25 + 25^2\times 35 + 35^2\times 20 + 45^2\times 5}{100} - (22,5)^2 = \frac{375+5625+21875+24500+10125}{100} - 506,25 = \frac{62500}{100} - 506,25 = 625 - 506,25 = 118,75\). Écart-type : \(\sigma = \sqrt{118,75} \approx 10,9\). Interprétation : les montants sont assez dispersés autour de la moyenne.
\end{enumerate}